% Part: computability
% Chapter: computability-theory
% Section: introduction

\documentclass[../../../include/open-logic-section]{subfiles}

\begin{document}

\olfileid{cmp}{thy}{int}
\olsection{مقدمة}

موضوع \emph{نظرية قابلية الحساب}، وهي فرع من المنطق، ما يتصل
بحساب الدوال والمجموعات، وحساب بعضها بالنسبة إلى بعض. وقد كان
هذا العلم يعرف بـ\emph{نظرية العودية}، وفي ذلك أثر نفوذ كليني في نشأته؛
ولا يزال الاسمان مستعملين.

الدالة~$f\colon\Nat\pto\Nat$ \emph{جزئية قابلة للحساب} متى تأتى
حسابها بأحد نماذج الحساب؛ فإن كانت $f$ كلية اقتصرنا على القول إن
$f$ \emph{قابلة للحساب}. وكذلك العلاقة~$R$ توصف بقابلية الحساب إذا كانت دالتها
المميزة~$\Char{R}$ قابلة للحساب. وإذا كانت $f$ و~$g$ دالتين
جزئيتين، كتبنا $f(x)\fdefined$ لنعني أن $f$ معرفة عند~$x$، أي إن
$x$ في مجال~$f$؛ وكتبنا $f(x)\fundefined$ لنعني العكس، أي إن $f$
غير معرفة عند~$x$. وأما قولنا $f(x)\simeq g(x)$ فمعناه أن $f(x)$ و
$g(x)$ إما غير معرفتين معًا، وإما معرفتان ومتساويتان معًا.

ولنا أن نبحث في قابلية الحساب من غير اختصاص بنموذج
معين؛ وذلك بإثبات دالة جزئية شاملة قابلة للحساب،
هي~$\fn{Un}(k,x)$. فبها يتأتى تعداد الدوال الجزئية القابلة للحساب.
ونجعل الرمز~$\cfind{k}$ للدالة الجزئية الأحادية الحجة
القابلة للحساب ذات الفهرس $k$، فيكون تعريفها
$\cfind{k}(x)\simeq\fn{Un}(k,x)$. (استعمل كليني $\{k\}$ لهذا الغرض،
لكن هذا الترميز قل استعماله في الآونة الأخيرة.) ولا يختص ذلك بالأحادية؛
بل نعدد، على وجه موحد، الدوال الجزئية القابلة للحساب بمختلف
رتبها، ونجعل $\cfind{k}[n]$ للدلالة على الدالة العودية الجزئية
ذات $n$ حجة وذات الفهرس $k$ في التعداد.

إذا كانت $f(\vec x,y)$ دالة كلية أو جزئية، فإن
$\umin{y}{f(\vec x,y)}$ هي الدالة في~$\vec x$ التي تعيد أصغر~$y$
بحيث $f(\vec x,y)=0$، على افتراض أن $f(\vec x,0)$، …،
$f(\vec x,y-1)$ كلها معرفة؛ وإذا لم يوجد $y$ كهذا، كانت
$\umin{y}{f(\vec x,y)}$ غير معرفة. وإذا كانت $R(\vec x,y)$ علاقة،
عرفنا $\umin{y}{R(\vec x,y)}$ بأنه أصغر~$y$ بحيث تصدق
$R(\vec x,y)$؛ أي أصغر~$y$ بحيث $1\tsub\Char{R}(\vec x,y)=0$.

ولإثبات قابلية دالة للحساب سبيلان:
\begin{enumerate}
\item السبيل الصارم: تعيين آلة تورنغ أو دالة عودية جزئية تعيينًا صريحًا،
  ثم إثبات أنها تحسب الدالة المطلوبة؛
\item السبيل غير الصوري: وصف خوارزمية للحساب، والاستناد إلى أطروحة تشرش.
\end{enumerate}
وليس بين السبيلين فصل حاد؛ فإن الوصف المفصل للخوارزمية ينبغي أن يبيّن
من أمرها ما يكفي لمعرفة كيف تُنشأ، من حيث المبدأ، آلة
تورنغ المناسبة، أو سلسلة التعريفات العودية الجزئية المناسبة. وقد تحجب
التفاصيل الرسمية الكاملة وجه الفائدة؛ فلنسلك طريقًا وسطًا بين
النهجين، فنورد براهين واضحة للحدس مستوفية للحجة، تدل على كيفية استخراج
التعريفات الدقيقة.

\end{document}
